% letter-keeptogether-closing-two-page.tex — layout stress: the closing block
% forced across the page-fill boundary on purpose (issue #421).
%
% The body is sized so that its last line lands close enough to the bottom of
% page one that \MakeCDossierClosing's two pieces — the closing text and the
% signature name — cannot both fit under it. Without a keep-together at that
% boundary the page builder takes the only legal break inside the block, at the
% \vspace{\CDossierLetterSignatureGapSkip}, leaving "Sincerely," alone at the
% foot of page one and "Ada Lovelace" alone at the head of page two. With it,
% the whole block moves to page two together and the break falls at the
% \addvspace above it instead.
%
% This is the letter's counterpart to the resume's *-keeptogether* fixtures:
% \clubpenalty and \widowpenalty cannot reach this boundary, because the closing
% and the name are two separate one-line paragraphs rather than the first and
% last line of one (#171, #342).
%
% Two independent checks cover it, and both were confirmed to fail with the
% \nobreak removed from \MakeCDossierClosing:
%
%   - the KEEPTOGETHER directive below, which requires the two texts on one
%     page. The signature name also appears in the page-two running header, so
%     the directive is decisive here because it compares the *last* page
%     carrying each text: split, the closing stays on page one while the name
%     is read from page two.
%   - the CLUB/WIDOW detector in page-break-check.awk, which sees the same split
%     geometrically. "Sincerely," ends mid-sentence, which is what makes a
%     one-line paragraph at a page edge reportable there.
\documentclass{careerdossier-letter}
% KEEPTOGETHER: Sincerely, ||| Ada Lovelace
\CDossierSetup{
  name     = {Ada Lovelace},
  headline = {Data Scientist},
  email    = {ada.lovelace@example.com},
  phone    = {+1 416 555 0142},
  location = {Toronto, Ontario, Canada},
}
\CDossierLetterSetup{
  date                   = {14 July 2026},
  recipient-name         = {Hiring Committee},
  recipient-title        = {Department of Applied Analytics},
  recipient-organization = {Northbridge Research Institute},
  recipient-address      = {123 Discovery Avenue\\Vancouver, British Columbia V6T 1Z4},
  subject                = {Application for the Senior Data Scientist position},
  salutation             = {Dear Members of the Hiring Committee,},
  closing                = {Sincerely,},
}
\begin{document}
\MakeCDossierLetterhead

I am writing to apply for the Senior Data Scientist position in computational
analytics at the Northbridge Research Institute. I recently completed a decade
of work spanning public-sector analytics, reporting platforms, and reproducible
research pipelines, and I am eager to contribute that experience to your group's
mission of building reliable, accessible scientific software.

My most recent work focused on developing scalable, reproducible reporting
services for high-volume public programs. I designed governed data pipelines,
implemented performance-critical processing routines, and collaborated closely
with program teams to translate operational requirements into maintainable
systems with measurable service standards. This work reduced duplicated
maintenance across teams, shortened incident investigation, and established
automated validation for data quality, reproducibility, and archival output.

Earlier in my career I built data-import services and validation tools for
public reporting datasets in a range of interchange formats. I improved
processing performance through batching, indexing, and query optimization while
preserving auditability, and I created internal libraries for schema validation,
error reporting, and automated documentation generation. Across these roles I
have consistently favoured tested, documented, and reproducible tools over
one-off analyses, because they are what allow a team to move quickly without
sacrificing reliability.

The position is especially appealing because of its emphasis on interdisciplinary
collaboration and dependable scientific software. My experience translating
analytical models into tested computational tools would allow me to contribute
quickly to your current projects, and I would welcome the opportunity to mentor
junior researchers and to support the group's documentation and reproducibility
practices.

Thank you for considering my application. I have enclosed my curriculum vitae and
would be pleased to discuss how my background and interests align with the goals
of your group. I am available at your convenience for a conversation and can
provide references, code samples, and further documentation on request. I would
be glad to walk through any of the projects described above in as much technical
detail as would be useful to the committee.

\MakeCDossierClosing
\end{document}
